\documentclass[12pt]{article}
\usepackage{amsmath}
\usepackage{amssymb}
\usepackage{geometry}
\usepackage{enumitem}
\geometry{margin=1in}

\setlist[itemize]{topsep=4pt,itemsep=2pt}
\setlist[enumerate]{topsep=4pt,itemsep=2pt}

\begin{document}

\section*{CSE400 – Fundamentals of Probability in Computing}

\subsection*{Lecture 10: Randomized Min-Cut Algorithm}

\bigskip

\section*{1. Outline (as given in slides)}

\begin{itemize}
\item Min-Cut Problem
\begin{itemize}
\item Why use min-cut?
\item What is min-cut?
\item Successful min-cut run
\item Unsuccessful min-cut run
\item Max-Flow Min-Cut Theorem
\end{itemize}

\item Deterministic Min-Cut Algorithm
\begin{itemize}
\item Stoer–Wagner minimum cut algorithm
\item Pseudocode
\end{itemize}

\item Randomized Min-Cut Algorithm
\begin{itemize}
\item Why randomized algorithm?
\item Karger’s Randomized Algorithm
\item Pseudocode
\item Comparison: Deterministic Min-Cut vs Randomized Min-Cut Algorithm
\item Theorem for min-cut set
\item Python Simulation
\item Class Participation
\end{itemize}
\end{itemize}

\bigskip

\section*{2. Min-Cut Problem}

\subsection*{2.1 Why Use Min-Cut?}

We use the min-cut algorithm in various applications to solve problems related to:

\begin{itemize}
\item Network connectivity
\item Reliability
\item Optimization
\end{itemize}

\subsubsection*{Applications(explicitly listed)}

\begin{enumerate}
\item \textbf{Network Design}
\begin{itemize}
\item Min cut helps in improving the efficiency of communication.
\item Optimizes network flow.
\item Used to find the minimum capacity cut.
\end{itemize}

\item \textbf{Communication Networks}
\begin{itemize}
\item Useful for understanding the vulnerability of networks to failures.
\item Helps in building robust and fault-tolerant communication networks.
\end{itemize}

\item \textbf{VLSI Design}
\begin{itemize}
\item Useful for partitioning circuits into smaller components.
\item Leads to reduced interconnectivity complexity.
\end{itemize}
\end{enumerate}

\bigskip

\subsection*{2.2 What is Min-Cut?}

\subsubsection*{Definition: Cut-Set}

A \textbf{cut-set} in a graph is:

\begin{quote}
A set of edges whose removal breaks the graph into two or more connected components.
\end{quote}

\bigskip

\subsubsection*{Definition: Minimum Cut}

Given a graph

[
G = (V,E)
]

with (n) vertices,

The \textbf{minimum cut (min-cut)} problem is:

\begin{quote}
To find a minimum cardinality cut-set in (G).
\end{quote}

Slides additionally state:

\begin{itemize}
\item Min-cut algorithms such as Karger’s algorithm are random.
\item They can be sensitive to the initial choice of edges.
\item If the algorithm contracts critical edges early, it might find a smaller cut.
\end{itemize}

\bigskip

\subsection*{2.3 Edge Contraction}

The main operation in the algorithm is \textbf{edge contraction}.

\subsubsection*{Definition}

Edge contraction is:

\begin{quote}
An operation that removes an edge from a graph while simultaneously merging the two vertices.
\end{quote}

\bigskip

\subsubsection*{Contracting an Edge ((u,v))}

When contracting edge ((u,v)):

\begin{enumerate}
\item Merge the two vertices (u) and (v) into one vertex.
\item Eliminate all edges connecting (u) and (v).
\item Retain all other edges in the graph.
\item The new graph:
\begin{itemize}
\item May have parallel edges
\item Has no self-loops
\end{itemize}
\end{enumerate}

The diagrams on pages 15–16 illustrate:

\begin{itemize}
\item Selection of an edge
\item Merging of its endpoints
\item Elimination of self-loops
\item Possible creation of parallel edges
\end{itemize}

\bigskip

\section*{3. Successful and Unsuccessful Min-Cut Runs}

\subsection*{3.1 Successful Min-Cut Run}

A \textbf{successful min-cut run}:

\begin{quote}
Refers to the success in the outcome of an algorithm designed to find the minimum cut in a graph.
\end{quote}

The figure on page 19 shows:

\begin{itemize}
\item A sequence of contractions
\item Final partition that corresponds to the minimum cut
\end{itemize}

\bigskip

\subsection*{3.2 Unsuccessful Min-Cut Run}

An \textbf{unsuccessful min-cut run}:

\begin{quote}
Refers to an iteration of a min-cut algorithm where the algorithm fails to correctly identify the minimum cut of a given graph.
\end{quote}

The figure on page 22 shows:

\begin{itemize}
\item Contraction of critical edges
\item Final partition that is not the true minimum cut
\end{itemize}

\bigskip

\section*{4. Max-Flow Min-Cut Theorem}

\subsection*{Statement}

\begin{quote}
In a flow network, the maximum amount of flow passing from the source to the sink is equal to the total weight of the edges in a minimum cut.
\end{quote}

\bigskip

\subsection*{Definitions(as introduced)}

\subsubsection*{Capacity of a Cut}

[
\text{Capacity of a cut} = \sum \text{capacity of edges in the cut oriented from a vertex} \in X \text{ to a vertex} \in Y
]

\bigskip

\subsubsection*{Minimum Cut}

\begin{itemize}
\item The cut in the network that has the smallest possible capacity.
\item \textbf{Minimum cut capacity} = the capacity of the minimum cut.
\end{itemize}

\bigskip

\subsubsection*{Maximum Flow}

\begin{itemize}
\item The largest possible flow from source (S) to sink (T).
\end{itemize}

The diagram on page 26 visually demonstrates:

\begin{itemize}
\item Partition of vertices
\item Flow values on edges
\item Equality between maximum flow and minimum cut capacity
\end{itemize}

\bigskip

\section*{5. Deterministic Min-Cut Algorithm}

\subsection*{5.1 Stoer–Wagner Min-Cut Algorithm}

Let (s) and (t) be two vertices of a graph (G).

Let (G/{s,t}) be the graph obtained by merging (s) and (t).

\subsubsection*{Theorem (exact wording from slides)}

A minimum cut of (G) can be obtained by taking the smaller of:

\begin{enumerate}
\item A minimum (s) - (t) - cut of (G), and
\item A minimum cut of (G/{s,t}).
\end{enumerate}

\bigskip

\subsubsection*{Justification}

The theorem holds since:

\begin{itemize}
\item Either there is a minimum cut of (G) that separates (s) and (t),\
$\rightarrow$ then a minimum (s) - (t) - cut of (G) is a minimum cut of (G);

\item Or there is none,\
$\rightarrow$ then a minimum cut of (G/{s,t}) does the job.
\end{itemize}

\bigskip

\subsection*{5.2 Pseudocode}

\subsubsection*{Algorithm 1: MinimumCutPhase(G, a)}

\begin{verbatim}
A ← {a};
while A ≠ V do
add to A the most tightly connected vertex;
return the cut weight as the "cut of the phase";
\end{verbatim}

\bigskip

\subsubsection*{Algorithm 2: MinimumCut(G)}

\begin{verbatim}
while |V| ≥ 1 do
choose any a from V;
MinimumCutPhase(G, a);
if the cut-of-the-phase is lighter than the current minimum cut then
store the cut-of-the-phase as the current minimum cut;
shrink G by merging the two vertices added last;
return the minimum cut;
\end{verbatim}

\bigskip

\subsection*{5.3 Time Complexity (Deterministic)}

Stoer–Wagner algorithm:

[
O(V \cdot E + V^2 \log V)
]

where:

\begin{itemize}
\item (V) = number of vertices
\item (E) = number of edges
\end{itemize}

\bigskip

\section*{6. Randomized Min-Cut Algorithm}

\subsection*{6.1 Why Randomized Algorithm?}

Randomized algorithms:

\begin{itemize}
\item Provide a probabilistic guarantee of success.
\item Provide a more accurate estimate of the minimum cut with fewer iterations.
\end{itemize}

Additional characteristics listed:

\begin{itemize}
\item Efficiency
\item Parallelization
\item Approximation Guarantees
\item Avoidance of Worst-Case Instances
\item Heuristic Nature
\item Robustness
\end{itemize}

\bigskip

\subsection*{6.2 Karger’s Randomized Algorithm}

The diagram on page 34 shows:

\begin{enumerate}
\item Pick edge (b), remove it and fuse its two corners into one.
\item Pick edge (d), remove it and fuse its two corners into one.
\end{enumerate}

The slide explicitly states:

\begin{quote}
The output cut produced here is {a,c,e}, but the minimal cut is either {b,e} or {a,d}.
\end{quote}

This example illustrates that random contraction can yield a non-minimal cut.

\bigskip

\subsection*{6.3 Pseudocode}

\subsubsection*{Algorithm 3: RECURSIVE-RANDOMIZED-MIN-CUT(G, $\alpha$)}

\textbf{Input:}

\begin{itemize}
\item Undirected multigraph (G) with (n) vertices
\item Integer constant ( \alpha > 0 )
\end{itemize}

\textbf{Output:}

\begin{itemize}
\item A cut (C) of (G)
\end{itemize}

\begin{verbatim}
if n ≤ α^3 then
C ← a min-cut of G found using brute force (exhaustive search);
else
for i ← 1 to α do
G' ← multigraph obtained by applying
n − ⌈n/√α⌉ random contraction steps on G;
C' ← RECURSIVE-RANDOMIZED-MIN-CUT(G', α);
if i = 1 or |C'| < |C| then
C ← C';
return C;
\end{verbatim}

\bigskip

\subsection*{6.4 Time Complexity (Randomized)}

Karger’s algorithm:

[
O(V^2)
]

where (V) is the number of vertices.

\bigskip

\section*{7. Comparison: Deterministic vs Randomized}

The slides state:

\begin{quote}
Any specific problem is responsible for deciding the approach.
\end{quote}

\bigskip

\subsection*{Deterministic Min Cut}

\begin{itemize}
\item Always guarantees an exact minimum cut.
\item May have higher time complexity for large graphs.
\item Time complexity:
[
O(V \cdot E + V^2 \log V)
]
\end{itemize}

\bigskip

\subsection*{Randomized Min Cut}

\begin{itemize}
\item Provides an approximate minimum cut with high probability.
\item Time complexity:
[
O(V^2)
]
\end{itemize}

\bigskip

\section*{8. Theorem for Min-Cut Set}

The algorithm outputs a min-cut set with probability at least:

[
\frac{2}{n(n-1)}
]

This is the probability guarantee stated in the slides.

\bigskip

\section*{9. Python Simulation}

Instructions given:

\begin{itemize}
\item Open Campuswire Post regarding today’s lecture (L10).
\item Download the `.ipynb` file pasted in the comments.
\end{itemize}

\bigskip

\section*{End of Lecture 10 Scribe}

\end{document}
